% Section 4 -- formal definitions. Definitional only: no figures, no macros.
% The "this is set cardinality, not structure" objection is raised here and
% answered here, which is cheaper than defending it after a reviewer raises it.

\section{Structural coverage: formal definitions}
\label{sec:definitions}

Let $G$ be the explicit subgraph of the mathlib dependency graph, as
characterized in Section~\ref{sec:is-explicit}: its nodes are declarations and
its edges are the source-attributable dependencies between them. A proof $p$
selects a set of premises $P(p) \subseteq V(G)$, the declarations it cites, and
those premises span an induced subgraph $G[P(p)]$. For a corpus $C$ we write
$V(C) = \bigcup_{p \in C} P(p)$ for its \emph{vocabulary}: the node set of the
subgraph the corpus's whole proof population induces.

\paragraph{Coverage.} Given a human corpus $H$ and a machine corpus $M$,
coverage is the family of quantities over $V(H)$ and $V(M)$: the two
vocabulary sizes and their ratio; the Jaccard index
$J = |V(H) \cap V(M)| \,/\, |V(H) \cup V(M)|$; and the human-only fraction
$|V(H) \setminus V(M)| \,/\, |V(H)|$. These quantities read no edges.
Computing $J$ consults only membership, so coverage is the node-set family and
breadth, below, the edge-derived one. Appendix~\ref{app:construct} states what
the graph does and does not supply to the coverage arithmetic, why we chose a
measure that needs no aggregation, and how each quantity is reported.

\paragraph{Breadth.} Structural breadth is the topological complement: the
number of distinct communities of a fixed partition of $G$ that a premise set
touches, optionally per unit of proof length. This does read edges, since the
partition is computed from them.

\paragraph{Which definitions need the hub-layer split.} Coverage does not. It
aggregates nothing, so Finding~3 does not bind, and the layer composition of
each vocabulary is reported in Section~\ref{sec:layer-composition} as a result
in its own right rather than as a correction. Breadth and every
centrality-flavoured quantity do need it. Breadth carries a second caveat, the
partition skew of Appendix~\ref{app:breadth}, and it is reported as secondary
throughout, with its three design-dependent magnitudes quoted together, never
one alone.
